\notechapter{N-20230209}{Vector Triple Product, Tensor Intuition, Linear Maps, and Matrix Representations}

\section{What the page is trying to prove}

The original handwritten pages are not a polished tensor note.  They begin with a vector identity, try to prove it geometrically and algebraically, then move toward tensors, linear maps, matrices, symmetry, and traces.  The thread is: vector operations have coordinate formulas, but the formulas are only shadows of coordinate-free objects.

The first main identity is the vector triple product
\[
  u\times(v\times w)=(u\cdot w)v-(u\cdot v)w.
\]
The note is trying to avoid treating this as a determinant trick only.  It asks why the right-hand side has to be a linear combination of $v$ and $w$, and then determines the coefficients.

\section{Geometric proof of the vector triple product}

Since $v\times w$ is perpendicular to the plane spanned by $v$ and $w$, the vector
\[
  u\times(v\times w)
\]
is perpendicular to $v\times w$.  Therefore it lies in the plane spanned by $v$ and $w$.  Hence there exist scalars $A,B$ such that
\[
  u\times(v\times w)=Av+Bw.
\]
This is the conceptual step: the expression has no reason to contain any direction outside the $v,w$-plane.

To determine $A$ and $B$, use scalar triple products.  Dot both sides with a vector perpendicular in the right way, or compare how both sides behave after dotting with $v$ and $w$.  The result is
\[
  A=u\cdot w,
  \qquad
  B=-(u\cdot v).
\]
Therefore
\[
  u\times(v\times w)=(u\cdot w)v-(u\cdot v)w.
\]
The identity is often remembered by the mnemonic “BAC minus CAB”, but the proof shows that it is really a projection formula inside the plane generated by $v$ and $w$.

\section{A determinant check}

One can also verify the identity in coordinates.  Write
\[
  u=(u_1,u_2,u_3),\quad v=(v_1,v_2,v_3),\quad w=(w_1,w_2,w_3).
\]
Then use
\[
  v\times w=
  \begin{pmatrix}
  v_2w_3-v_3w_2\\
  v_3w_1-v_1w_3\\
  v_1w_2-v_2w_1
  \end{pmatrix}.
\]
Expanding $u\times(v\times w)$ and collecting the coefficients of $v$ and $w$ gives the same expression.  The coordinate proof is longer, but it is a useful check that no orientation sign has been lost.

The note's calculation seems to move back and forth between geometric and coordinate reasoning.  That is good mathematical instinct: the geometry tells us what shape the answer must have, and coordinates verify the signs.

\section{Scalar triple product and volume}

The scalar triple product is
\[
  [u,v,w]=(u\times v)\cdot w.
\]
It is the signed volume of the parallelepiped spanned by $u,v,w$.  Equivalently,
\[
  [u,v,w]=\det(u,v,w).
\]
It is alternating: swapping two vectors changes the sign; if two vectors are equal or linearly dependent, the volume is zero.

This explains why identities involving cross and dot products often reduce to determinant identities.  The determinant is the coordinate expression of oriented volume.

\section{From vectors to tensors}

The later pages begin to use tensor language.  A vector is not its coordinate column.  A vector is the geometric object; a coordinate column is how it is represented after choosing a basis.  If the basis changes, the coordinate column changes.

A covector, or linear functional, transforms in the dual way.  If vectors transform by a matrix $A$, covectors transform by the inverse transpose.  A tensor is a multilinear object with several vector and covector slots.  Its components transform by the corresponding combination of $A$, $A^{-1}$, $A^T$, or $A^{-T}$.

This is the motivation behind the note's scattered formulas involving transposes and inverse matrices.  They are not arbitrary matrix manipulations; they are bookkeeping rules for preserving the same object under coordinate changes.

\section{Linear maps and matrices}

A linear map $T:V\to W$ satisfies
\[
  T(av+bw)=aT(v)+bT(w).
\]
After choosing a basis of $V$ and a basis of $W$, the map is represented by a matrix.  But the same map has a different matrix in a different basis.  If the basis of $V$ changes by $P$ and the basis of $W$ changes by $Q$, the representing matrix changes by a conjugation-like rule involving $Q^{-1}$ and $P$.

For an endomorphism $T:V\to V$, changing basis gives
\[
  [T]_{new}=P^{-1}[T]_{old}P.
\]
Thus similar matrices represent the same linear operator in different bases.  This is why trace, determinant, eigenvalues, and characteristic polynomial are meaningful: they are invariant under similarity.

\section{Symmetric, skew-symmetric, and trace ideas}

The handwritten pages also touch on symmetric matrices.  A matrix $A$ is symmetric if
\[
  A^T=A,
\]
and skew-symmetric if
\[
  A^T=-A.
\]
Over complex vector spaces, the corresponding notion is Hermitian:
\[
  A^*=A.
\]
The trace is
\[
  \operatorname{tr}(A)=\sum_i A_{ii}.
\]
It satisfies
\[
  \operatorname{tr}(AB)=\operatorname{tr}(BA),
\]
which implies trace is invariant under similarity:
\[
  \operatorname{tr}(P^{-1}AP)=\operatorname{tr}(A).
\]
This is a clean example of the general theme: the matrix entries change, but the underlying linear map has invariants.


\section{Basis change and invariant meaning}

If a linear operator has matrix \(A\) in one basis and the basis changes by \(P\), then the new matrix is
\[
  P^{-1}AP.
\]
The entries may change completely, but trace, determinant, characteristic polynomial, and eigenvalues remain the same.  These quantities belong to the operator, not to the chosen coordinate description.

This is the central lesson behind the tensor discussion: coordinate arrays are useful only when we know how they transform.

\section{Expanded synthesis and advanced context}

The note begins with $u\times(v\times w)$ and ends near tensors.  That is a natural progression.  The vector triple product says a coordinate-looking identity is really a geometric projection formula.  The scalar triple product says determinants measure oriented volume.  Tensor language says coordinate arrays must transform correctly if they are to represent genuine objects.  Linear algebra becomes mature when one stops identifying an object with one of its coordinate descriptions.



\paragraph{Expanded synthesis.}
For this chapter, the elementary material should be read as a study of what remains invariant when coordinates change. The earlier sections (`What the page is trying to prove`, `Geometric proof of the vector triple product`, `A determinant check`, `Scalar triple product and volume`, `From vectors to tensors`) compute with arrays, bases, pivots, determinants, or eigenvectors; the deeper object is a linear map together with the spaces it acts on. A matrix is a representation of that map after bases have been chosen. This is why formulas such as
\[
  A' = P^{-1}AP,\qquad \widetilde A = T^{-1}AS
\]
are not cosmetic: they distinguish an intrinsic morphism from a coordinate description.

The structural hierarchy is as follows. Kernels record what a map forgets, images record what it reaches, cokernels record obstructions to solvability, and quotients record information after an equivalence has been imposed. Determinants act on the top exterior power and therefore measure oriented volume; traces are invariant under cyclic rearrangement and become characters in representation theory; adjoints depend on an inner product and lead to self-adjoint spectral theory.

A useful way to return to the computations is to ask, for every formula, which invariant it is measuring. Row reduction measures rank and kernel; diagonalization measures decomposition into invariant subspaces; Gram--Schmidt measures orthogonal projection; positive definiteness measures ellipsoid geometry; tensor notation measures how components transform. The elementary methods are therefore the finite-dimensional laboratory in which duality, quotienting, functoriality, and spectral decomposition can be seen clearly.


\section{Elementary-to-advanced bridge: linear maps versus matrices}

A linear map $T:V\to W$ is independent of coordinates.  A matrix appears only after choosing a basis of $V$ and a basis of $W$.  If $[v]_B$ is the coordinate column of $v$ in basis $B$, then a matrix $A$ represents $T$ by
\[
  [T(v)]_C=A[v]_B.
\]
Changing bases changes the matrix, not the map.

For an endomorphism $T:V\to V$, if $P$ is the transition matrix from a new basis to an old one, the matrix changes by similarity:
\[
  A_{new}=P^{-1}A_{old}P.
\]
Therefore trace, determinant, rank, and characteristic polynomial are invariants of the linear map, because they survive similarity.

\section{Elementary-to-advanced bridge: diagrams for basis change}

Basis change is easiest to remember as a commuting diagram:
\[
  \begin{array}{ccc}
  V&\xrightarrow{T}&V\\
  \downarrow [\,]_B&&\downarrow [\,]_B\\
  k^n&\xrightarrow{A_B}&k^n
  \end{array}
\]
Changing the coordinate isomorphism $[\,]_B$ changes the bottom arrow by conjugation.  The actual top arrow is unchanged.

This is a functorial idea: choosing coordinates is an isomorphism from an abstract vector space to $k^n$.  Matrix algebra is the transported version of coordinate-free linear algebra.

\section{Elementary-to-advanced bridge: bilinear maps and tensor products}

A bilinear map
\[
  B:V\times W\to U
\]
can be represented as a linear map out of a tensor product:
\[
  \widetilde B:V\otimes W\to U,\qquad \widetilde B(v\otimes w)=B(v,w).
\]
This universal property explains why tensor products appear whenever two linear inputs interact.

Matrix multiplication itself is a contraction of tensors: indices repeated once upstairs and once downstairs are summed.  The elementary array notation becomes systematic tensor calculus once the variance of indices is tracked.

\section{Returning to the note}

The original correction from vector triple product to linear maps and matrices is mathematically natural.  Vector identities teach invariant geometric operations; matrix representations teach how invariant maps acquire coordinate arrays.  The bridge between them is tensor language: it records which parts are geometric and which parts are basis-dependent.

% Second-pass deepening for waves 2-4: wave 2
\section{Second-pass deepening: basis change as naturality of coordinates}

The note's basis-change formulas become more meaningful when written as isomorphisms.  A basis $B$ is an isomorphism
\[
  [\,]_B:V\to k^n.
\]
A linear map $T:V\to W$ has coordinate representation
\[
  [T]_{C,B}=[\,]_C\circ T\circ [\,]_B^{-1}.
\]
Changing bases means replacing the coordinate isomorphisms, not changing $T$.

This formula immediately explains why maps between different spaces transform by
\[
  A' = S^{-1}AT
\]
while endomorphisms transform by similarity
\[
  A'=P^{-1}AP.
\]
The first has two independent coordinate changes; the second uses the same space on both sides.

\section{Second-pass deepening: tensors as multilinear natural transformations}

A tensor of type $(r,s)$ on a vector space $V$ can be viewed as a multilinear map with $r$ vector outputs and $s$ covector inputs, or equivalently as an element of
\[
  V^{\otimes r}\otimes (V^*)^{\otimes s}.
\]
Basis change rules are not arbitrary; they are forced by requiring the underlying multilinear object to be independent of coordinates.

The original transition from vector identities to matrices becomes a lesson in invariance.  Coordinate arrays are allowed to change, but the geometric or algebraic object must not.
